% GENERATED FILE. Edit canonical lessons, then run npm run generate:books.
% canonical-source-hash: fnv1a64:ae0bd39a3e29d93a
% canonical-lessons: ML-C561-kacceri, ML-C561-vadyam, ML-C561-kottuka-drum, ML-C561-girrar, ML-C561-otakkulal

\chapter{Drums, Guitars and Flutes}
\label{ch:ml-drums-guitars-and-flutes}
\hlchaptermodality{561}

\begin{chapteropening}
\textbf{By the end of this chapter:} \emph{I can say a concert, a musical instrument, to beat a drum, a guitar and a flute.}

Complete the last lesson of chapter 561: I can say a concert, a musical instrument, to beat a drum, a guitar and a flute.
\end{chapteropening}

\section[\texorpdfstring{\ml{കച്ചേരി} (kaccēri)}{കച്ചേരി (kaccēri)}]{\ml{കച്ചേരി} (kaccēri) — a concert; a recital}

\label{lesson:ML-C561-kacceri}



\begin{quote}
Before the new one: say the Malayalam for a sculptor, then the Malayalam for an exhibition.
\end{quote}

\subsection*{You'll want to know: \ml{കച്ചേരി}}
\textbf{\ml{കച്ചേരി}} — \emph{kaccēri} — \textquotedblleft{}a concert; a recital\textquotedblright{}.

\begin{quote}
\ml{ഇന്നലെ} \ml{രാത്രി} \ml{ഞങ്ങൾ} \ml{ഒരു} \ml{സംഗീത} \ml{കച്ചേരിക്ക്} \ml{പോയി}.
\end{quote}

\begin{quote}
\emph{innale rātri ñaṅṅaḷ oru saṁgīta kaccērikkŭ pōyi.}
\end{quote}

\begin{quote}
Last night we went to a music concert.
\end{quote}

A classical concert; the same word also names a court or a government office.

\begin{cousinweb}
\textbf{\ml{കച്ചേരി}} (\emph{kaccēri}) is usually traced to the North Indian word for a court-house, Hindi \emph{kacaharī}, which CDIAL sets under a reconstructed compound meaning court-house (CDIAL 3429). Malayalam keeps the sense of a court or an office; how it came to mean a concert is not settled.
\end{cousinweb}

\begin{grammarlens}[title={What you've built}]
One more word for describing what you lived and felt.
\end{grammarlens}

\subsection*{Your turn}
\begin{itemize}
\raggedright
  \item \emph{Say it:} \emph{kaccēri}
  \item \emph{Say it:} \emph{kaccēri}, once more
  \item \emph{Say it:} \emph{innale rātri ñaṅṅaḷ oru saṁgīta kaccērikkŭ pōyi.} — Last night we went to a music concert.
  \item \emph{Recall:} say the Malayalam for a sculptor, then the Malayalam for an exhibition, then say \emph{kaccēri} again
\end{itemize}

\begin{tcolorbox}[breakable,colback=teal!4,colframe=teal!35!black,title={Before you move on}]
What does \ml{കച്ചേരി} mean? (\textquotedblleft{}A concert; a recital\textquotedblright{}.) Say it once more.
\end{tcolorbox}
\section[\texorpdfstring{\ml{വാദ്യം} (vādyaṁ)}{വാദ്യം (vādyaṁ)}]{\ml{വാദ്യം} (vādyaṁ) — a musical instrument}

\label{lesson:ML-C561-vadyam}



\begin{quote}
Before the new one: say the Malayalam for an exhibition, then the Malayalam for a concert.
\end{quote}

\subsection*{You'll want to know: \ml{വാദ്യം}}
\textbf{\ml{വാദ്യം}} — \emph{vādyaṁ} — \textquotedblleft{}a musical instrument\textquotedblright{}.

\begin{quote}
\ml{ഉത്സവത്തിന്} \ml{പല} \ml{വാദ്യങ്ങൾ} \ml{ഉണ്ടായിരുന്നു}.
\end{quote}

\begin{quote}
\emph{utsavattinŭ pala vādyaṅṅaḷ uṇṭāyirunnu.}
\end{quote}

\begin{quote}
There were many instruments at the festival.
\end{quote}

Plural: \ml{വാദ്യങ്ങൾ}, \emph{vādyaṅṅaḷ}.

\ml{വാദ്യോപകരണം} is the full formal word for a musical instrument.

\begin{cousinweb}
\textbf{\ml{വാദ്യം}} (\emph{vādyaṁ}) is borrowed from Sanskrit \emph{vādya}, \textquotedblleft{}to be played\textquotedblright{}, and as a noun \textquotedblleft{}instrumental music\textquotedblright{} (CDIAL 11511).
\end{cousinweb}

\begin{grammarlens}[title={What you've built}]
One more word for describing what you lived and felt.
\end{grammarlens}

\subsection*{Your turn}
\begin{itemize}
\raggedright
  \item \emph{Say it:} \emph{vādyaṁ}
  \item \emph{Say it:} \emph{vādyaṁ}, once more
  \item \emph{Say it:} \emph{utsavattinŭ pala vādyaṅṅaḷ uṇṭāyirunnu.} — There were many instruments at the festival.
  \item \emph{Recall:} say the Malayalam for an exhibition, then the Malayalam for a concert, then say \emph{vādyaṁ} again
\end{itemize}

\begin{tcolorbox}[breakable,colback=teal!4,colframe=teal!35!black,title={Before you move on}]
What does \ml{വാദ്യം} mean? (\textquotedblleft{}A musical instrument\textquotedblright{}.) Say it once more.
\end{tcolorbox}
\section[\texorpdfstring{\ml{കൊട്ടുക} (koṭṭuka)}{കൊട്ടുക (koṭṭuka)}]{\ml{കൊട്ടുക} (koṭṭuka) — to beat a drum}

\label{lesson:ML-C561-kottuka-drum}



\begin{quote}
Before the new one: say the Malayalam for a concert, then the Malayalam for a musical instrument.
\end{quote}

\subsection*{You'll want to know: \ml{കൊട്ടുക}}
\textbf{\ml{കൊട്ടുക}} — \emph{koṭṭuka} — \textquotedblleft{}to beat a drum\textquotedblright{}.

\begin{quote}
\ml{ഉത്സവത്തിന്} \ml{അവർ} \ml{ചെണ്ട} \ml{കൊട്ടി}.
\end{quote}

\begin{quote}
\emph{utsavattinŭ avar ceṇṭa koṭṭi.}
\end{quote}

\begin{quote}
They played the chenda drum at the festival.
\end{quote}

Past: \ml{കൊട്ടി}, \emph{koṭṭi}.

Also to knock or tap; \ml{കൈ} \ml{കൊട്ടുക} is to clap.

\begin{cousinweb}
\textbf{\ml{കൊട്ടുക}} (\emph{koṭṭuka}) is native Dravidian: DEDR 2063 gives it as to beat so as to make a sound, as on a drum or a bell (DEDR 2063).
\end{cousinweb}

\begin{grammarlens}[title={What you've built}]
One more word for describing what you lived and felt.
\end{grammarlens}

\subsection*{Your turn}
\begin{itemize}
\raggedright
  \item \emph{Say it:} \emph{koṭṭuka}
  \item \emph{Say it:} \emph{koṭṭuka}, once more
  \item \emph{Say it:} \emph{utsavattinŭ avar ceṇṭa koṭṭi.} — They played the chenda drum at the festival.
  \item \emph{Recall:} say the Malayalam for a concert, then the Malayalam for a musical instrument, then say \emph{koṭṭuka} again
\end{itemize}

\begin{tcolorbox}[breakable,colback=teal!4,colframe=teal!35!black,title={Before you move on}]
What does \ml{കൊട്ടുക} mean? (\textquotedblleft{}To beat a drum\textquotedblright{}.) Say it once more.
\end{tcolorbox}
\section[\texorpdfstring{\ml{ഗിറ്റാർ} (giṟṟār)}{ഗിറ്റാർ (giṟṟār)}]{\ml{ഗിറ്റാർ} (giṟṟār) — a guitar}

\label{lesson:ML-C561-girrar}



\begin{quote}
Before the new one: say the Malayalam for a musical instrument, then the Malayalam for to beat a drum.
\end{quote}

\subsection*{You'll want to know: \ml{ഗിറ്റാർ}}
\textbf{\ml{ഗിറ്റാർ}} — \emph{giṟṟār} — \textquotedblleft{}a guitar\textquotedblright{}.

\begin{quote}
\ml{എന്റെ} \ml{സുഹൃത്ത്} \ml{ഗിറ്റാർ} \ml{നന്നായി} \ml{വായിക്കും}.
\end{quote}

\begin{quote}
\emph{enṟe suhṛttŭ giṟṟār nannāyi vāyikkuṁ.}
\end{quote}

\begin{quote}
My friend plays the guitar well.
\end{quote}

One plays an instrument with \ml{വായിക്കുക}, the same verb as to read.

\begin{cousinweb}
\textbf{\ml{ഗിറ്റാർ}} (\emph{giṟṟār}) was borrowed from English \emph{guitar}.
\end{cousinweb}

\begin{grammarlens}[title={What you've built}]
One more word for describing what you lived and felt.
\end{grammarlens}

\subsection*{Your turn}
\begin{itemize}
\raggedright
  \item \emph{Say it:} \emph{giṟṟār}
  \item \emph{Say it:} \emph{giṟṟār}, once more
  \item \emph{Say it:} \emph{enṟe suhṛttŭ giṟṟār nannāyi vāyikkuṁ.} — My friend plays the guitar well.
  \item \emph{Recall:} say the Malayalam for a musical instrument, then the Malayalam for to beat a drum, then say \emph{giṟṟār} again
\end{itemize}

\begin{tcolorbox}[breakable,colback=teal!4,colframe=teal!35!black,title={Before you move on}]
What does \ml{ഗിറ്റാർ} mean? (\textquotedblleft{}A guitar\textquotedblright{}.) Say it once more.
\end{tcolorbox}
\section[\texorpdfstring{\ml{ഓടക്കുഴൽ} (ōṭakkuḻal)}{ഓടക്കുഴൽ (ōṭakkuḻal)}]{\ml{ഓടക്കുഴൽ} (ōṭakkuḻal) — a flute}

\label{lesson:ML-C561-otakkulal}



\begin{quote}
Before the new one: say the Malayalam for to beat a drum, then the Malayalam for a guitar.
\end{quote}

\subsection*{You'll want to know: \ml{ഓടക്കുഴൽ}}
\textbf{\ml{ഓടക്കുഴൽ}} — \emph{ōṭakkuḻal} — \textquotedblleft{}a flute\textquotedblright{}.

\begin{quote}
\ml{ഓടക്കുഴലിന്റെ} \ml{ശബ്ദം} \ml{ദൂരെ} \ml{നിന്ന്} \ml{കേട്ടു}.
\end{quote}

\begin{quote}
\emph{ōṭakkuḻalinṟe śabdaṁ dūre ninnŭ kēṭṭu.}
\end{quote}

\begin{quote}
The sound of a flute was heard from far away.
\end{quote}

\ml{ഓട} a kind of reed + \ml{കുഴൽ} a tube.

\begin{cousinweb}
\textbf{\ml{ഓടക്കുഴൽ}} (\emph{ōṭakkuḻal}) is native Dravidian, a compound of \ml{ഓട}, \textquotedblleft{}a large reed\textquotedblright{} (DEDR 1043), and \ml{കുഴൽ}, \textquotedblleft{}a tube, a flute\textquotedblright{} (DEDR 1818).
\end{cousinweb}

\begin{grammarlens}[title={What you've built}]
One more word for describing what you lived and felt.
\end{grammarlens}

\subsection*{Your turn}
\begin{itemize}
\raggedright
  \item \emph{Say it:} \emph{ōṭakkuḻal}
  \item \emph{Say it:} \emph{ōṭakkuḻal}, once more
  \item \emph{Say it:} \emph{ōṭakkuḻalinṟe śabdaṁ dūre ninnŭ kēṭṭu.} — The sound of a flute was heard from far away.
  \item \emph{Recall:} say the Malayalam for to beat a drum, then the Malayalam for a guitar, then say \emph{ōṭakkuḻal} again
\end{itemize}

\begin{tcolorbox}[breakable,colback=teal!4,colframe=teal!35!black,title={Before you move on}]
What does \ml{ഓടക്കുഴൽ} mean? (\textquotedblleft{}A flute\textquotedblright{}.) Say it once more.
\end{tcolorbox}
